
\section{Conclusiones}
El proyecto permitió representar el problema que trae la selección de materiales para mejorar la eficiencia energetica de una vivienda creando un modelo de optimización. Una de las dificultades principales fue lograr retratar esta situacion real utilizando variables para acercarnos lo mas posible a la realidad. Los supuestos utilizados son importantes ya que sin ellos la dificultad de la optimización a realizar incrementa al integrarse variables las cuales varian de manera muy difícil de predecir.

Respecto a la generación de instancias factibles, hubo gran dificultad ya que no bastaba con generar valores aleatorios, estos habían de ser definidos en un rango y verificar que ciertas inecuaciones se cumplan. Para asegurar la existencia de instancias viables el generador determina un escenario de máxima eficiencia térmica y utiliza este resultado para establecer un valor de $U^{max}$ que permita mantener la factibilidad de la instancia. De esta forma, las 15 instancias generadas cumplen con la condición de factibilidad definida para el proyecto.

La informática nos permitió la automatización del proceso de generación de las instancias, pudimos implementar el proceso entero mediante múltiples funciones y diccionarios en conjunto a listas, las cuales facilitan establecer valores mínimos y máximos para variables de manera rápida. Esto, junto a la utilización de una semilla, permite generar instancias de manera aleatoria pero manteniendo la posibilidad de reproducir los mismos resultados. Además, la separación del proceso en distintas funciones facilita modificar los parámetros del generador y adaptar la generación de instancias según el tamaño requerido. Gracias a esto entre otras cosas el proyecto tiene buena capacidad de reproducción y escalabilidad.

Finalmente, para generar instancias de mayor tamaño o reducir los tiempos de ejecución, se podrían incorporar estructuras de datos más eficientes, optimizar el proceso de validación y evitar cálculos repetidos durante la generación. También sería posible paralelizar la generación de instancias independientes, permitiendo aprovechar varios núcleos del procesador. Estas mejoras permitirían aumentar la cantidad y tamaño de las instancias sin incrementar proporcionalmente el tiempo requerido para generarlas. Sobre el uso de recursos del algoritmo, en este caso no se estan realizando muchas instancias por lo que no ha de haber un gran consumo de recursos o tiempo, si se quisiera expandir la cantidad de instancias, posibles optimizaciones podría ser implementar el generador en un lenguaje como C++, que puede ofrecer un mejor rendimiento en determinadas operaciones, aunque esto implicaría una mayor complejidad de implementación.
